% Resume
\thispagestyle{empty}

\chapter*{Resum\'e}
\addcontentsline{toc}{chapter}{Resum\'e}
\markboth{Resum\'e}{Resum\'e}

\section{Motivácia práce}

Diplomová práca sa venuje využitiu dôkazov s nulovou znalosťou v ekosystéme
Ethereum, konkrétne návrhu post-kvantovo bezpečného ZkVM založeného na
mriežkovej kryptografii. Hlavnou motiváciou práce je problém škálovateľnosti
Etherea. V súčasnom modeli musia uzly siete pri spracovaní blokov opätovne
vykonávať a overovať všetky transakcie. Tento prístup je síce bezpečný, ale pri
väčšom zaťažení siete vedie k vysokým výpočtovým nárokom a obmedzuje celkovú
priepustnosť systému.

Jedným z hlavných smerov, ako tento problém riešiť, je použitie dôkazov s
nulovou znalosťou. Tie umožňujú dokázať správnosť výpočtu bez toho, aby musel
overovateľ celý výpočet vykonať znova. V prostredí Ethereum sa preto objavuje
myšlienka tzv. snarkifikácie, teda nahradenia priameho overovania vykonania
blokov overením stručného kryptografického dôkazu. Ak je možné správnosť
vykonania bloku potvrdiť jediným dôkazom, otvára sa cesta k efektívnejšiemu
overovaniu, jednoduchšej škálovateľnosti a do budúcnosti aj k prirodzenejšej
integrácii rollupových riešení. Širší kontext Ethereum a snarkifikácie je
spracovaný v kapitole~\ref{chap:analysis}, konkrétne v sekciách
\ref{sec:ethereum} a~\ref{sec:snarkification}. Vývoj od ZkEVM k všeobecným
ZkVM, ako aj úloha rekurzie a skladania, sú opísané v sekcii~\ref{sec:zkevm}
a znázornené na obrázkoch~\ref{fig:ivc_diagram} a~\ref{fig:ivc-folding}.

Popri škálovateľnosti je dôležitou motiváciou aj bezpečnosť. Mnohé súčasné ZKP
systémy využívajú predpoklady, ktoré nie sú odolné voči kvantovým útokom.
Časť ekosystému síce smeruje k post-kvantovým riešeniam založených na hašoch, napríklad k
systémom typu STARK, no tie majú často menej priaznivé vlastnosti z pohľadu
veľkosti dôkazu alebo nákladov na verifikáciu. Práve preto sa v práci skúma
mriežková kryptografia ako alternatíva, ktorá ponúka post-kvantovú bezpečnosť
založenú na problémoch ako SVP a SIS a zároveň zachováva algebraickú štruktúru
vhodnú pre konštrukciu pokročilých dôkazových systémov.

Mriežková kryptografia je pre túto prácu dôležitá aj z architektonického
hľadiska. Nie je zaujímavá iba ako náhrada za klasické schémy založené na
eliptických krivkách, ale aj ako základ pre skladanie dôkazov v štýle IVC.
Práce ako Ajtaijove záväzky, LaBRADOR, Greyhound a najmä LatticeFold ukazujú,
že mriežkové primitíva možno využiť aj pri návrhu stručných a rekurzívnych
dôkazových mechanizmov. Práca preto skúma, či možno tieto prístupy použiť na
vytvorenie všeobecného RISC-V ZkVM, ktorá by bola použiteľná aj pre úlohy
súvisiace s dokazovaním správneho vykonania Ethereum blokov. Teoretické
vlastnosti mriežkovej kryptografie sú rozobraté v sekcii~\ref{sec:lattices} a
výber použitých kryptografických primitív v návrhu riešenia v sekcii
~\ref{sec:crypto}.

Ďalšou motiváciou práce je výber vhodného výpočtového modelu. Priame dokazovanie
EVM je zložité, pretože EVM nebolo navrhnuté s ohľadom na efektívnu ZK
aritmetizáciu. Obsahuje komplikovanejšie inštrukcie, pracuje s 256-bitovými
slovami a používa model, ktorý je pre tvorbu dôkazov menej prirodzený. Preto sa
v súčasnosti pozornosť presúva k všeobecným ZkVM nad architektúrou RISC-V,
ktorá je jednoduchšia, registrová a vhodnejšia na formálne zachytenie vykonania
programu. V tejto práci je RISC-V zvolený ako vhodný základ na návrh systému,
ktorý by spájal post-kvantovú bezpečnosť, architektúru ZkVM a potenciálne
využitie v prostredí Ethereum. Porovnanie vlastností existujúcich ZkVM je
uvedené v tabuľke~\ref{tab:zkvm-comparison}.

Motiváciou práce je teda spojiť tri hlavné požiadavky: zlepšenie
škálovateľnosti, zachovanie post-kvantovej bezpečnosti a návrh realistickej
architektúry všeobecného ZkVM. Z tejto motivácie prirodzene vyplýva aj samotný
problém, ktorému sa práca venuje.

\section{Vymedzenie problému}

Na základe analýzy súčasného stavu možno povedať, že oblasť ZkVM sa rýchlo
rozvíja, no zároveň je veľmi rozmanitá. Existujú riešenia založené na rôznych
inštrukčných sadách, rôznych modeloch aritmetizácie aj na rôznych dôkazových
mechanizmoch. Napriek tomu však v literatúre ani v dostupnom open-source
ekosystéme zatiaľ neexistuje praktické riešenie všeobecného RISC-V ZkVM, ktoré
by bolo postavené na mriežkových primitívach a využívalo skladanie ako základ
inkrementálne overiteľného výpočtu.

Práve táto medzera tvorí hlavný problém diplomovej práce. Hoci mriežková
kryptografia a schémy ako LatticeFold teoreticky naznačujú, že post-kvantové
skladanie dôkazov je možné, nie je zrejmé, či sa tieto myšlienky dajú pretaviť
do funkčného systému pre všeobecné vykonávanie programov. Nestačí totiž iba
navrhnúť jednotlivé kryptografické primitíva. Je potrebné vytvoriť celú
architektúru, v ktorej bude možné reprezentovať stav virtuálneho stroja,
aritmetizovať jednotlivé kroky vykonania, kontrolovať konzistenciu pamäte a
spájať všetky kroky do jedného IVC reťazca. Tento problém je formulovaný v
kapitole~\ref{chap:problem_statement}. Základná schéma navrhovanej architektúry
je rozpracovaná v kapitole návrhu riešenia~\ref{chap:design}, najmä v sekciách
\ref{sec:arch}, \ref{sec:riscv_circuit} a~\ref{sec:ivc}, a znázornená na
obrázkoch~\ref{fig:execution_flow} a~\ref{fig:snark_wrap}.

Ďalšou časťou problému je výkonnosť takéhoto riešenia. To, že je nejaký prístup
teoreticky možný, ešte neznamená, že bude aj prakticky použiteľný. Pri návrhu
ZkVM založeného na mriežkovej kryptografii vznikajú náklady spojené s
aritmetizáciou inštrukcií, generovaním dôkazových dát, skladaním inštancií a
rekurzívnym overovaním. Bez konkrétnej implementácie nie je možné reálne
posúdiť, ako sa takýto systém správa, aké sú jeho časové a pamäťové nároky a
ako sa porovnáva s existujúcimi ZkVM založenými na klasických alebo hash-based
post-kvantových prístupoch.

Dôležitou súčasťou problému je aj otázka dlhých výpočtov. Ak má byť takýto
systém zaujímavý pre Ethereum, nestačí ukázať, že funguje na malých ukážkových
príkladoch. Musí byť aspoň na prototypovej úrovni overené, či je schopný zvládať
dlhšie výpočtové stopy bez zjavného zlyhania. V prípade mriežkovej kryptografie
je táto otázka obzvlášť dôležitá aj kvôli rastu šumu a praktickým obmedzeniam,
ktoré sa môžu objaviť pri dlhšom rekurzívnom skladaní.

Problém riešený v tejto diplomovej práci možno preto zhrnúť nasledovne: chýba
konkrétna a experimentálne overená inštancia všeobecného RISC-V ZkVM založeného
na mriežkovej kryptografii, ktorá by umožnila posúdiť jeho realizovateľnosť,
výkonnostné kompromisy a vhodnosť pre dlhšie výpočty. Práca na túto medzeru
reaguje návrhom a implementáciou prototypu. Cieľom nie je vytvoriť okamžite
produkčne použiteľný systém, ale overiť, či je takýto smer z architektonického
a implementačného hľadiska uskutočniteľný a aké obmedzenia sa pri ňom prejavia.

\section{Návrh riešenia}

Návrh riešenia je opísaný v kapitole~\ref{chap:design}. Jeho cieľom je vytvoriť
architektúru všeobecného RISC-V ZkVM, ktorá bude vhodná pre dokazovanie dlhých
výpočtov a zároveň bude postavená na mriežkových kryptografických primitívach.
Navrhovaný systém je rozdelený na niekoľko vrstiev. Na vstupnej strane vchádza
program napísaný vo vyššom programovacom jazyku, typicky v Ruste, ktorý sa
skompiluje do binárnej podoby pre 32-bitovú architektúru RISC-V. Tento binárny
súbor následne vykonáva virtuálny stroj, pričom počas vykonávania vzniká stopa
obsahujúca informácie o stave stroja, dekódovanej inštrukcii a prípadných
pamäťových operáciách. Celkový tok spracovania je znázornený na
obrázku~\ref{fig:execution_flow}.

Základným stavebným prvkom navrhovaného riešenia je tzv. krokový obvod $F$,
ktorý je opísaný v sekcii~\ref{sec:riscv_circuit}. Tento obvod reprezentuje
jeden krok vykonania RISC-V programu. Jeho úlohou je overiť, že prechod medzi
stavom pred vykonaním inštrukcie a stavom po jej vykonaní zodpovedá pravidlám
zvolenej ISA. To zahŕňa načítanie inštrukcie, jej dekódovanie, kontrolu
správnosti operandov, vykonanie samotnej semantiky inštrukcie a aktualizáciu
stavu stroja. Zvolený kryptografický postup pracuje nad algebraickou
štruktúrou z Goldilocks cyklotomického poľa, bolo potrebné
aritmetizovať 32-bitovú logiku procesora do formy vhodnej pre CCS. Práve CCS
predstavuje model obmedzení, na ktorom je návrh postavený.

Samotný obvod $F$ však nestačí na dokazovanie dlhých výpočtov. Z tohto dôvodu
je v sekcii~\ref{sec:ivc} definovaný rozšírený obvod $F'$, ktorý slúži ako
jadro IVC mechanizmu. Tento obvod obaluje jeden krok vykonania a zároveň
zabezpečuje integritu celého reťazca dôkazov. V každom kroku kontroluje, že
aktuálny stav nadväzuje na predchádzajúci stav, že predchádzajúci akumulátor je
platný a že predošlé skladanie prebehlo korektne. Na tento účel je do návrhu
zaradený aj vnorený overovateľ pre LatticeFold. Výsledkom je architektúra,
v ktorej sa jednotlivé kroky výpočtu neoverujú izolovane, ale sú postupne
spájané do jedného priebežne aktualizovaného akumulátora.

Z kryptografického hľadiska je návrh riešenia založený na kombinácii viacerých
stavebných blokov, ktoré sú rozobraté v sekcii~\ref{sec:crypto}. Základom je
model IVC, ktorý umožňuje dokazovať správnosť dlhého výpočtu inkrementálne.
Praktickú realizáciu tejto myšlienky zabezpečuje skladanie inštancií, inšpirované
prácami Nova a HyperNova, no v tejto práci je klasický mechanizmus nahradený
schémou LatticeFold. Tá poskytuje post-kvantové bezpečnostné vlastnosti a
zároveň umožňuje budovať rekurzívny dôkazový mechanizmus nad mriežkovými
predpokladmi. Dôležitou súčasťou návrhu je aj použitie Poseidon2 záväzkov na
reťazenie krokov a identifikáciu stavu IVC, ako aj voľba obaľovacieho SNARKu, kde sa
uvažuje o spojení SuperSpartan a Greyhound, opísané v
sekcii~\ref{subsec:wrapping_snark}. Záverečný krok obalenia výsledného artefaktu
je znázornený na obrázku~\ref{fig:snark_wrap}.

Osobitnou otázkou návrhu bola práca s pamäťou. Priamo vkladať úplnú kontrolu
pamäte do každého kroku výpočtu by znamenalo výrazné zvýšenie zložitosti
obvodu. Preto návrh využíva odloženú kontrolu konzistencie pamäte, opísanú v
sekcii~\ref{subsec:memory_check}. Počas vykonávania sa zaznamenávajú všetky
čítania a zápisy do pamäťového logu, pričom ich úplná kontrola sa presúva až do
neskoršej fázy. Tento prístup znižuje záťaž krokového obvodu a zároveň zachováva
možnosť overiť správnosť read-over-write správania pamäte.

Z pohľadu architektúry je dôležité, že návrh riešenia nie je iba zoznamom
samostatných kryptografických techník, ale tvorí prepojený systém. Sekcia
\ref{sec:arch} opisuje vysokú úroveň architektúry, sekcia
\ref{sec:riscv_circuit} formalizuje samotný krok vykonania a sekcia
\ref{sec:ivc} vysvetľuje, ako sa tieto kroky spájajú do dlhého dôkazového
reťazca. Cieľom návrhu nie je deklarovať finálne optimalizované riešenie, ale
ukázať, že je možné vytvoriť konzistentnú architektúru, v ktorej sa spoja
RISC-V vykonávanie, CCS aritmetizácia, IVC, skladanie pomocou LatticeFold a
post-kvantový obaľovací mechanizmus do jedného uceleného ZkVM návrhu.

\section{Implementácia}

Implementácia navrhnutého riešenia je opísaná v kapitole~\ref{chap:implementation}.
Jej cieľom bolo vytvoriť funkčný prototyp, na ktorom bude možné overiť praktickú
uskutočniteľnosť navrhnutej architektúry. Implementácia nepokrýva všetky časti,
ktoré by boli potrebné pre plnohodnotný produkčný ZkVM, ale zahŕňa hlavné
stavebné bloky potrebné na experimentálne overenie konceptu. Súčasťou
prototypu je podmnožina inštrukcií architektúry \texttt{rv32imac}, CCS obmedzenia
pre ich vykonávanie, CCS obmedzenia pre Poseidon2 záväzky a taktiež
obmedzenia potrebné pre vnútro obvodového overovateľa.

Prototyp bol implementovaný v jazyku Rust a použitý technologický základ je
opísaný v sekcii~\ref{sec:tech_stack}. Algebraická a hašovacia vrstva je
postavená na knižniciach z ekosystému Plonky3, samotné skladanie využíva
implementáciu LatticeFold od Nethermind. Ako hosťovský program bol v
rámci experimentov použitý aj \texttt{REVM}, čo umožnilo skúmať správanie
systému na dlhších výpočtoch súvisiacich s Ethereum prostredím.

Významnou súčasťou implementácie bolo vytvorenie vlastného RISC-V
emulátora, ktorý generuje vykonávaciu stopu potrebnú pre dôkazový systém.
Emulátor vykonáva program po jednotlivých inštrukciách a pri každom kroku
sprístupňuje úplný stav stroja. Tým sa oddelila samotná exekúcia programu od
generovania dôkazových dát. Vykonávacia stopa sa následne používa pri tvorbe
dôkazového vektora a pri napĺňaní CCS štruktúr. Vysokoúrovňový pohľad na
implementovaný tok vykonania je znázornený na obrázku~\ref{fig:execution_flow_overview}.

Z implementačného hľadiska patrili medzi najnáročnejšie časti integrácia
kryptografických knižníc a syntéza obmedzení. Bežné kryptografické knižnice sú
navrhnuté tak, aby vracali najmä finálny výsledok operácie, zatiaľ čo v ZkVM je
potrebné zachytiť aj medzistavy výpočtu na účely witness generovania. Preto bola
vytvorená vlastná implementácia permutácie Poseidon2, ktorá sprístupňuje vnútorné
stavy potrebné pre CCS obmedzenia. Podobne bolo potrebné upraviť aj prácu s
knižnicou LatticeFold, aby bolo možné znovu použiť interné verifierové a pomocné
rutiny pri konštrukcii in-circuit logiky.

Ďalšou dôležitou časťou implementácie bola syntéza CCS obmedzení. Tá zahŕňala
preklad logiky RISC-V inštrukcií do algebraickej formy, explicitné zachytenie
výpočtu Poseidon2 a aj konštrukciu seba referenčných obmedzení potrebných pre
overovateľa skladania. Keďže tieto časti závisia od konečnej štruktúry celého CCS,
implementácia musela prebiehať vo viacerých fázach. Súčasťou práce boli aj
praktické optimalizácie, napríklad zdieľanie maticových štruktúr medzi podobnými
typmi obmedzení, aby zostal prototyp spustiteľný aj na bežnom hardvéri.

Výsledkom implementácie nie je úplne dokončený produkčný systém, ale funkčný
prototyp hlavných architektonických komponentov. Tento prototyp umožnil overiť,
že navrhnutý prístup nie je iba teoretický, ale že jeho kľúčové časti možno
skutočne spojiť do jedného bežiaceho systému a následne experimentálne
vyhodnotiť.

\section{Vyhodnotenie}

Vyhodnotenie prototypu je spracované v kapitole~\ref{chap:evaluation}. Jeho
cieľom nie je preukázať produkčnú pripravenosť systému, ale posúdiť, či
navrhnutá architektúra funguje v praxi a aké výsledky dosahuje neúplný
prototyp.

Základný benchmark bol vykonaný na programe počítajúcom 100. Fibonacciho číslo.
Pri tomto behu vznikla vykonávacia stopa dlhá 16 inštrukcií a zaznamenali sa 2
pamäťové operácie. Celkový čas skladania dosiahol 688,24 sekundy. Už tento
jednoduchý príklad ukázal, že aktuálny prototyp je výrazne pomalší než moderné
ZkVM riešenia.

Podrobnejšie porovnanie s inými systémami je uvedené v tabuľke~\ref{tab:zkvm-benchmarks-v2}.
Na rovnakom hardvéri dosiahla implementácia tejto práce čas generovania dôkazu
11:28.92 pri maximálnej spotrebe 12.34 GB RAM. Pre porovnanie, Pico dosiahlo
2:04.50 a 6.27 GB RAM, SP1 dosiahlo 1:50.17 a 9.96 GB RAM a ZKsync Airbender
4:06.73 pri 24.34 GB RAM. Z týchto výsledkov vyplýva, že navrhovaný systém je v
aktuálnej podobe pomalší ako porovnávané riešenia. Jednou z možných budúcich
výhod však môže byť menšia veľkosť finálneho dôkazu.

Popri malom benchmarku bol vykonaný aj dlhší experiment s hosťovským programom
založeným na \texttt{revm}. Tento beh trval približne 9 hodín a 29 minút, počas
ktorých systém pokračoval vo vykonávaní programu a overovaní IVC reťazca.
Maximálna spotreba pamäte sa pohybovala okolo 13.7 GB. Hoci ani tento výsledok
neznamená pripravenosť na produkčné dokazovanie Ethereum blokov, je dôležitý tým,
že ukazuje schopnosť systému pokračovať v dlhšom výpočte bez zjavného kolapsu.

Vyhodnotenie sa venuje aj otázke, prečo je aktuálny prototyp pomalší. Hlavným
praktickým problémom sa ukázala byť algebraická náročnosť skladania, najmä práca
s multilineárnymi rozšíreniami a súvisiacimi operáciami v rámci LatticeFold.
Profilovanie ukázalo, že dominantná časť výpočtu neleží v samotnom vykonávaní
virtuálneho stroja, ale v algebraických operáciách potrebných pre skladanie.
Druhým možným dôvodom je samotná voľba CCS ako výpočtového modelu.

Na základe vykonaných experimentov práca odpovedá aj na tri výskumné otázky.
Na otázku realizovateľnosti mriežkového RISC-V ZkVM odpovedá čiastočne kladne:
na architektonickej úrovni bol takýto systém navrhnutý a na prototypovej úrovni
aj implementovaný, hoci zatiaľ neúplne. Na otázku výkonových kompromisov je
odpoveď v súčasnom stave skôr negatívna, keďže implementácia je pomalšia než
porovnávané ZkVM. Na otázku vhodnosti pre dlhé výpočty sú výsledky opatrne
pozitívne, pretože dlhší experiment neukázal okamžitý problém s nekontrolovaným
rastom pamäte ani zlyhanie rekurzívneho reťazca, hoci na definitívne závery by
bolo potrebné rozsiahlejšie testovanie.

Celkové vyhodnotenie teda ukazuje, že navrhovaný smer je architektonicky
uskutočniteľný, ale zatiaľ výkonovo nekonkurencieschopný voči etablovaným
riešeniam. Prototyp však potvrdil, že mriežková kryptografia sa dá použiť ako
základ pre všeobecný RISC-V ZkVM a že tento smer má zmysel ďalej rozvíjať.

\section{Záver}

Záver práce je spracovaný v kapitole~\ref{chap:conclusion}. Diplomová práca sa
venovala návrhu a prototypovej implementácii post-kvantovo bezpečného RISC-V
ZkVM založeného na mriežkovej kryptografii s potenciálnym využitím v prostredí
Ethereum. Hlavným prínosom práce je návrh ucelenej architektúry, ktorá spája
RISC-V vykonávanie, CCS aritmetizáciu, IVC, skladanie pomocou LatticeFold a
uvažovaný obaľovací mechanizmus založený na kombinácii SuperSpartan a
Greyhound.

Práca zároveň ukázala, že tento smer nie je iba teoretickou úvahou. Podarilo sa
implementovať funkčný prototyp hlavných stavebných blokov systému, ktorý
umožnil experimentálne overiť základnú realizovateľnosť navrhovaného prístupu.
Výsledky síce ukazujú, že aktuálna implementácia zaostáva za etablovanými ZkVM
z pohľadu výkonu a nemožno ju považovať za produkčne pripravené riešenie, no
zároveň naznačujú, že hlavné obmedzenia sú skôr implementačné a aritmetizačné
ako principiálne.

Z pohľadu výskumného prínosu je dôležité najmä to, že práca adresuje konkrétnu
medzeru v súčasnom stave poznania. Ukazuje, ako možno mriežkové primitíva
využiť pri návrhu všeobecného ZkVM a aké praktické problémy pri tom vznikajú.
Zároveň poskytuje prvé experimentálne údaje, na základe ktorých možno
realistickejšie hodnotiť výkonové kompromisy, vhodnosť zvoleného výpočtového
modelu a správanie systému pri dlhších výpočtoch.

Práca tým vytvára základ pre ďalší výskum. Prirodzenými pokračovaniami sú
efektívnejší výpočtový model než CCS, lepšia algebraická optimalizácia,
využitie SIMD alebo GPU akcelerácie a tiež prechod na novšie mriežkové schémy,
napríklad LatticeFold+ alebo Neo. Celkovo možno uzavrieť, že navrhnutý smer je
perspektívny a zaslúži si ďalšie rozpracovanie, aj keď v súčasnom stave ešte
nepredstavuje konkurencieschopné produkčné riešenie.
